%
% 31-warum.tex -- Warum Potenzreihen nicht immer reichen
%
% (c) 2026 Prof Dr Andreas Müller
%
\subsection{Warum Potenzreihen nicht immer reichen
\label{buch:differentialgleichungen:fuchs:subsection:warum}}
Löst man die Differentialgleichung
\eqref{buch:differentialgleichungen:fuchs:eqn:dgl}
nach $y''$ auf, erhält man
\[
y''
=
-
\frac{p(x)}{x}y'
-
\frac{q(x)}{x^2}y
\]
und kann erkennen, dass schon für konstante Funktionen $p(x)$ und
$q(x)$ die zweite Ableitung $y''$ an der Stelle $x=0$ eine
Singulartität haben muss, die sich mit einer gewöhnlichen
Potenzreihe nicht darstellen lässt.
Nimmt man eine physikalische Anwendung an, dann haben
$x^2y''$, $xy'$ und $y$ alle die gleiche Masseinheit, so dass
die Funktionen $p(x)$ und $q(x)$ dimensionslos sein können.

%
% Konstante Koeffizienten
%
\subsubsection{Konstante Koeffizienten
\label{buch:differentialgleichungen:fuchs:subsection:konstant}}
Noch etwas expliziter kann man versuchen, die Differentialgleichung
mit $p(x)=p_0$ und $q(x)=q_0$ mit einem Ansatz
\[
y(x) = \sum_{k=0}a_kx^k
\]
zu lösen.
Setzt man den Ansatz in die Gleichung ein, erhält man
\begin{align*}
x^2
\sum_{k=0}^\infty k(k-1) a_k x^{k-2}
+
p_0x\sum_{k=0}^\infty ka_k x^{k-1}
+
q_0\sum_{k=0}^\infty a_kx^k
&=
0
\\
\sum_{k=0}^\infty k(k-1) a_k x^k
+
p_0\sum_{k=0}^\infty ka_k x^k
+
q_0\sum_{k=0}^\infty a_kx^k
&=
0
\\
\sum_{k=0}^\infty
\bigl(
k(k-1) + p_0k + q_0
\bigr)
a_k
x^k
&=
0.
\end{align*}
Damit die Reihe verschwindet, müssen alle Koeffizienten verschwinden,
es gilt also entweder 
\[
k(k-1)+p_0k+q_0
=
k^2+(p_0-1)k+q_0
=
0
\qquad\text{oder}\qquad
a_k=0
\]
für alle $k\ge 0$ gelten.
Für $k=0$ besagt die Gleichung, dass $q_0=0$ oder $a_0=0$ sein
muss.
Für $k=1$ folgt, dass $p_0-1+q_0=0$ oder $a_1=0$ sein muss.
Auf der linken Seite steht eine quadratische Gleichung für $k$, diese 
Bedingung kann also höchstens für zwei Koeffizienten erfüllt sein,
entsprechend können höchstens zwei der Koeffizienten der Potenzreihe
$y(x)$ von $0$ verschieden sein.
Sind $k_1$ und $k_2$ die zugehörigen Koeffizienten, dann muss
\[
k_1k_2 = q_0
\qquad\text{und}\qquad
k_1+k_2
=
1-p_0
\]
sein.
Ausser für diese speziellen Werte von $p_0$ und $q_0$ ist
die quadratische Gleichung nie erfüllt, entsprechend müssen
alle Koeffizienten $a_k=0$ sein und damit ist $y(x)=0$
die triviale Lösung.

Der Lösungsansatz mit einer Potenzreihe führt also nicht zu einer
nichttrivialen Lösung.

%
% Lösungsversuch mit einer verallgemeinerten Potenzreihe
%
\subsubsection{Lösungsversuch mit einer verallgemeinerten Potenzreihe}
Wir versuchen die Differentialgleichung 
\eqref{buch:differentialgleichungen:fuchs:eqn:dgl}
mit konstanten Koeffizienten $p(x)=p_0$ und $q(x)=q_0$
mithilfe einer verallgemeinerten Potenzreihe
\begin{align*}
y(x)
&=
\sum_{k=0}^\infty
a_k x^{k+\varrho}
\intertext{zu lösen.
Die Ableitungsterme sind}
p_0xy'(x)
&=
\sum_{k=0}^\infty
p_0
(k+\varrho)
a_k x^{k+\varrho}
\\
x^2y''
&=
\sum_{k=0}^\infty
(k+\varrho)(k+\varrho-1)a_kx^{k+\varrho}.
\end{align*}
Setzt man diese in die Differentialgleichung ein, entsteht die Gleichung
\begin{align*}
\sum_{k=0}^\infty
(k+\varrho)(k+\varrho-1)
a_kx^{k+\varrho}.
+
\sum_{k=0}^\infty
p_0
(k+\varrho)
a_k x^{k+\varrho}
+
\sum_{k=0}^\infty
q_0
a_k x^{k+\varrho}
&=
0
\\
\sum_{k=0}^\infty
\bigl(
(k+\varrho)(k+\varrho-1)
+
p_0
(k+\varrho)
+
q_0
\bigr)
a_k x^{k+\varrho}
&=
0.
\end{align*}
Durch Koeffizientenvergleich folgt, dass
\[
\bigl((k+\varrho)^2+(p_0-1)(k+\varrho) + q_0\bigr) a_k = 0
\]
für alle $k$ sein muss.
Dies ist eine quadratische Gleichung, es 
folgt also wieder, dass höchstens zwei der Koeffizienten $a_k$ von $0$
verschieden sein können, die Lösung müsste also von der Form
\[
y(x)
=
A x^{\lambda_1} + B x^{\lambda_2}
\]
sein, wobei $\lambda_1$ und $\lambda_2$ Lösungen der quadratischen
Gleichung
\[
\lambda^2 +(p_0-1)\lambda + q_0
=
0
\]
ist.
Diese Lösung kann man natürlich auch direkt aus einem
Potenzansatz erhalten.
